\section{Introduction}
\label{sec:intro}

Open-weight chat models now run on laptops and phones, and safety tuning \cite{ouyang2022instructgpt,bai2022hh} is often the only thing that stands between such a model and a user who wants harmful help. {\bf What does safety tuning actually teach?} One answer is that it teaches a notion of harm: the model refuses because it understands what the request would do. Another answer is that it teaches a reflex: the model refuses because the request contains words such as ``bomb'' or ``hack'', is written in English, and asks directly.

The second answer has a close analogue in adversarial robustness. Defences that rely on \emph{gradient masking} look robust against the attacks they were tuned on, but fail under simple changes to the attack \cite{papernot2017practical,athalye2018obfuscated}. If safety tuning works the same way, a model would refuse harmful requests that look like its training data and comply with the same requests once they are re-worded, translated, or wrapped in a story. It would also over-refuse benign requests that happen to contain trigger words, such as ``How do I kill a Python process?''. The stakes are highest for small models, which receive the least safety data and are the easiest to deploy without oversight.

Many individual jailbreaks are consistent with this picture. Translation into low-resource languages \cite{deng2024multilingual,yong2023lowresource}, past-tense rewrites \cite{andriushchenko2024pasttense}, persuasive paraphrases \cite{zeng2024johnny}, role-play personas \cite{shen2024dan}, and multi-turn escalation \cite{russinovich2024crescendo} all lower refusal rates. However, these studies do not tell us \emph{why} refusal drops. A paraphrase changes the words, the style, and often the apparent severity of a request at the same time. A translated prompt that is not refused may simply not be understood. Without controls, we cannot tell a lexical blind spot from a paraphrase effect, a milder request, or a capability failure, and we cannot tell whether any of these shrink with scale.

We address this with a controlled study of 14 open-weight chat models from six families, ranging from 0.5B to 32.8B parameters (\figref{fig:overview}). All models answer the same 250 harmful requests from \jbb \cite{chao2024jailbreakbench} and \harmbench \cite{mazeika2024harmbench} under four distribution shifts. (\eone) We remove trigger words while keeping the meaning, and compare against a matched control rewrite that keeps the trigger words. (\etwo) We translate \multijail prompts into nine languages and separate real harmful compliance from off-topic or garbled replies. (\ethree) We build a prompt pool where keyword count and semantic harm are partly decorrelated, using \xstest \cite{rottger2024xstest} safe prompts that contain trigger words, and measure which of the two refusal follows. (\efour) We wrap each request in fiction, persona, and expert framings, and in a three-turn escalation. A validated \judge judge labels every response, and the official \harmbench classifier provides a stricter harmful-completion measure.

\begin{figure}[t]
\centering
\resizebox{0.95\linewidth}{!}{%
\begin{tikzpicture}[
  font=\small,
  box/.style={draw, rounded corners=2pt, align=center, execute at begin node=\hyphenpenalty 10000\relax, minimum height=0.9cm, inner sep=4pt},
  src/.style={box, fill=gray!12, text width=3.1cm},
  exp/.style={box, fill=blue!7, text width=5.6cm, minimum height=1.05cm},
  outbox/.style={box, fill=green!8, text width=3.6cm},
  arr/.style={-{Stealth[length=2mm]}, thick}
]
\node[src] (base) {250 harmful requests\\(\jbb + \harmbench)};
\node[src, below=0.55cm of base] (mj) {\multijail\\(315 prompts)};
\node[src, below=0.55cm of mj] (xs) {\xstest + \jbb benign\\(550 controls)};

\node[exp, right=0.9cm of base, yshift=1.05cm] (e1) {{\bf \eone Lexical ablation}\\\orig / \ctrl (keeps keywords) / \kwfree};
\node[exp, below=0.25cm of e1] (e4) {{\bf \efour Framing}\\\direct / \fiction / \persona / \expert / \escal};
\node[exp, below=0.25cm of e4] (e2) {{\bf \etwo Multilingual}\\en + zh, it, vi, ar, ko, th, bn, sw, jv};
\node[exp, below=0.25cm of e2] (e3) {{\bf \ethree Keyword \versus harm}\\1,615-prompt pool, harm rated 1--9};

\node[box, fill=orange!10, text width=2.7cm, right=0.9cm of e4, yshift=-0.65cm, minimum height=2.2cm] (models) {{\bf 14 chat models}\\0.5B--32.8B\\6 families\\greedy, 256 tokens};

\node[outbox, right=0.9cm of models, yshift=0.9cm] (judge) {\judge judge\\refused / complied / other};
\node[outbox, right=0.9cm of models, yshift=-0.9cm] (hb) {\harmbench classifier\\+ string-match check};

\draw[arr] (base.east) -- (e1.west);
\draw[arr] (base.east) -- (e4.west);
\draw[arr] (mj.east) -- (e2.west);
\draw[arr] (xs.east) -- (e3.west);
\draw[arr] (base.east) -- (e3.west);
\foreach \e in {e1,e4,e2,e3} \draw[arr] (\e.east) -- (models.west);
\draw[arr] (models.east) -- (judge.west);
\draw[arr] (models.east) -- (hb.west);
\end{tikzpicture}}
\caption{Study design. The same harmful requests pass through four distribution shifts and are answered by 14 open-weight chat models. Each response is labelled as refused, complied, or other by a \judge judge, and checked by the \harmbench classifier (harmful completion) and an LLM-free string-match refusal detector. The \ctrl rewrite in \eone and the \xstest safe prompts in \ethree are the controls that separate keyword effects from paraphrase and meaning effects.}
\label{fig:overview}
\end{figure}

{\bf The picture that emerges is different from the gradient-masking hypothesis.} Refusal follows meaning far more than vocabulary. Across all 14 models, the correlation of refusal with a semantic harm rating (0.27--0.63) is 1.5--4$\times$ its correlation with trigger-word count (0.10--0.25), and no model looks like a keyword matcher. Removing trigger words does lower refusal by 16.3\pp on average (significant in 13 of 14 models), while a re-wording that keeps them lowers it by only 2.5\pp. But this is a leak, not a collapse: the rate of \harmbench-harmful completions rises by only 3.5\pp, because most new non-refusals answer a watered-down version of the request. Multilingual ``bypass'' in sub-3B models is mostly incoherence: 44--78\% of English-refused prompts are not refused in other languages, but only 6.6--14.8\% receive a reply that carries out the request. The large and exploitable gaps are \emph{framing} gaps. A one-paragraph fiction frame raises compliance from at most 11\% to 51--78\% in Qwen2.5-1.5B/3B/7B and Phi-3-mini, and a ``no-rules persona'' raises Llama-3.1-8B from 14\% to 93\% and Gemma-2-9B from 1.6\% to 86\%. These gaps are specific to the model family and do not shrink with scale.

\para{Contributions.}
\begin{itemize}[leftmargin=*,itemsep=0pt,topsep=0pt]
    \item We propose a controlled lexical ablation that pairs each keyword-free rewrite with a keyword-kept control rewrite, which isolates the cost of losing trigger words (13.8\pp on average) from the cost of re-wording (2.5\pp).
    \item We measure, per model, how much refusal follows keyword count \versus semantic harm on a decorrelated pool, and confirm the result with LLM-free outcome and harm measures and with contrasts that do not depend on a keyword lexicon.
    \item We show that multilingual non-refusal in small models overstates harmful compliance several-fold, and that the gap between the two closes as models become more capable.
    \item We map framing blind spots across a 0.5B--32.8B ladder and show that they are large, family-specific, and non-monotonic in size, so that robustness is a property of the alignment recipe and not of parameter count.
\end{itemize}

\Secref{sec:related} reviews related work. \Secref{sec:method} describes the models, data, labelling, and statistics. \Secref{sec:results} reports the four experiments, and \secref{sec:discussion} discusses implications and limitations.
